\documentclass[10p,a4paper]{scrartcl}
%-------------------------------------------
%---Packages--------------------------------
%-------------------------------------------
\usepackage[utf8]{inputenc}
%\usepackage[T1]{fontenc}
%\usepackage{txfonts}
\usepackage{amsmath}
\usepackage{amsthm}
\usepackage{amsfonts}
\usepackage{array}
\usepackage{amssymb}
\usepackage{blindtext}
\usepackage{caption}
\usepackage{color}
\usepackage{csquotes}	    %
\usepackage{enumitem}	    %pour mieux bosser avec les listes. ajoute option label
\usepackage[yyyymmdd]{datetime}        %pour définir date custom
\usepackage{etaremune}
\usepackage{environ}
\usepackage{fancybox}
\usepackage{fancyhdr} 	    % Custom headers and footers
\usepackage{fancyref}
%\usepackage{float}
\usepackage{floatrow}       %float and floatrow can't be together...
\usepackage{gensymb}
\usepackage{graphicx}
\usepackage[colorlinks=true, linkcolor=purple, citecolor=cyan]{hyperref}
\usepackage{footnotebackref}
\usepackage{lipsum}
\usepackage{mathtools}
\usepackage{multicol}	    %gérer plusieurs colonnes
\usepackage{setspace}
\usepackage{subcaption}
\usepackage{todonotes}	    %Bonne gestion des TODOs
%TODO commenté pour tester l'utilité... à voir% \usepackage[tc]{titlepic}      %Permet de mettre une image en page de garde
\usepackage{tikz}	    % Pour outil de dessin puissant
\usepackage{ulem}	    %underline sur plusieurs lignes (avec \uline{})
\usepackage{vmargin} 	    %gestion des marges, avec dans l'ordre : gauche, haut, droit, bas, en-tête, entre en-tête et texte, bas de page, hauteur entre bas de page et texte
\usepackage{wrapfig}
\usepackage{xcolor}
\usepackage{xparse}                    %Pour utiliser NewDocumentCommand et des arguments 'mmooo'
%\usepackage{fullpage} 	    %supprime toutes les marges allouées aux notes, aussi en haut et en bas

%\ExplSyntaxOn
\pagestyle{fancyplain}	    %Makes all pages in the document conform to the custom headers and footers

%-------------------------------------------
%---Document Commands-----------------------
%---------------------------{----------------
\NewDocumentCommand{\framecolorbox}{oommm}
 {% #1 = width (optional)
  % #2 = inner alignment (optional)
  % #3 = frame color
  % #4 = background color
  % #5 = text
  \IfValueTF{#1}%
   {\IfValueTF{#2}%
    {\fcolorbox{#3}{#4}{\makebox[#1][#2]{#5}}}%
    {\fcolorbox{#3}{#4}{\makebox[#1]{#5}}}%
   }%
   {\fcolorbox{#3}{#4}{#5}}%
 }%
%------------------------------------------------
%------------------ENGLISH----------------------
%----------------------------------------------

\NewDocumentCommand{\epflTitle}{mO{Olivier Cloux}O{\today}O{Notes de Cours en}D<>{../../Common}}%Arguments : Matière, Auteur, Date, Titre du doc
{
\begin{titlepage}
    \vspace*{\fill}
    \begin{center}
        \normalfont \normalsize
        \textsc{Ecole Polytechnique Fédérale de Lausanne} \\ [25pt] % Your university, school and/or department name(s)
        \textsc{#4} %Titre du doc
        \\ [0.4 pt]
        \horrule{0.5pt} \\[0.4cm] % Thin top horizontal rule
        \huge #1 \\ % Matière
        \horrule{2pt} \\[0.5cm] % Thick bottom horizontal rule
        \includegraphics[width=8cm]{#5/EPFL_logo}
        ~\\[0.5 cm]
        \small\textsc{#2}\\[0.4cm]
        \small\textsc{#3}\\
        ~\\
        ~\\
        \includegraphics[scale=0.5]{#5/creativeCommons}
    \end{center}
    \vspace*{\fill}
\end{titlepage}
}


%-------------------------------------------
%-------------MATH NEW COMMANDS-------------
%-------------------------------------------
\newcommand{\somme}[2]{\ensuremath{\sum\limits_{#2}^{#1}}}
\newcommand{\produit}[2]{\ensuremath{\prod\limits_{#2}^{#1}}}
\newcommand{\limite}{\lim\limits_}
\newcommand{\llimite}[3]{\limite{\substack{#1 \\ #2}}\left(#3\right)}	%limites à deux condiitons
\newcommand{\et}{\mbox{ et }}
\newcommand{\deriv}[1]{\ensuremath{\, \mathrm d #1}}	%sigle dx, dt,dy... des dérivées/intégrales
%\newcommand{\fx}{\ensuremath{f'(\textbf{x}_0 + h}}
\newcommand{\ninf}{\ensuremath{n \to \infty}}	       %pour les limites : n tend vers l'infini
\newcommand{\xinf}{\ensuremath{x \to \infty}}	       %pour les limites : x tend vers l'infini
\newcommand{\infint}{\ensuremath{\int_{-\infty}^{\infty}}}
\newcommand{\xo}{\ensuremath{x \to 0}}									%x to 0
\newcommand{\no}{\ensuremath{n \to 0}}									%n zéro
\newcommand{\xx}{\ensuremath{x \to x}}									%x to x
\newcommand{\Xo}{\ensuremath{x_0}}										%x zéro
\newcommand{\X}{\ensuremath{\mathbf{X}} }
\newcommand{\A}{\ensuremath{\mathbf{A}} }
\newcommand{\R}{\ensuremath{\mathbb{R}} }								%ensemble de R
\newcommand{\rn}{\ensuremath{\mathbb{R}^n} } 							%ensemble de R de taille n
\newcommand{\Rm}{\ensuremath{\mathbb{R}^m} }  							%ensemble de R de taille m
\newcommand{\C}{\ensuremath{\mathbb{C}} }
\newcommand{\N}{\ensuremath{\mathbb{N}} }
\newcommand{\Z}{\ensuremath{\mathbb{Z}} }
\newcommand{\Q}{\ensuremath{\mathbb{Q}} }
\newcommand{\rtor}{\ensuremath{\R \to \R} }
\newcommand{\pour}{\mbox{ pour }}
\newcommand{\coss}[1]{\ensuremath{\cos\(#1\)}}						%cosinus avec des parenthèses de bonne taille (genre frac)
\newcommand{\sinn}[1]{\ensuremath{\sin\(#1\)}}					%sinus avec des parentèses de bonne taille (genre frac)
\newcommand{\txtfrac}[2]{\ensuremath{\frac{\text{#1}}{\text{#2}}}}		%Fractions composées de texte
\newcommand{\evalfrac}[3]{\ensuremath{\left.\frac{#1}{#2}\right|_{#3}}}
\renewcommand{\(}{\left(}												%Parenthèse gauche de taille adaptive
\renewcommand{\)}{\right)}
\newcommand{\longeq}{=\joinrel=}												%Parenthèse droite de taille adaptive


%-------------------------------------------------------
%------------------MISC NEW COMMANDS--------------------
%-------------------------------------------------------
\newcommand{\degre}{\ensuremath{^\circ}}
%\newdateformat{\eudate}{\THEYEAR-\twodigit{\THEMONTH}-\twodigit{\THEDAY}}



%-------------------------------------------------------
%------------------TEXT NEW COMMANDS--------------------
%-------------------------------------------------------
\newcommand{\ts}{\textsuperscript}
\newcommand{\evid}[1]{\textbf{\uline{#1}}}        %mise en évidence (gras + souligné)



%\newcommand{\Exemple}{\underline{Exemple}}
\newcommand{\Theoreme}{\underline{Théorème}}
\newcommand{\Remarque}{\underline{Remarque}}
\newcommand{\Definition}{\underline{Définition} }
\newcommand{\skinf}{\sum^{\infty}_{k=0}}
\newcommand{\combi}[2]{\ensuremath{\begin{pmatrix} #1 \\ #2 \end{pmatrix}}}	%combinaison parmi 1 de 2
\newcommand{\intx}[3]{\ensuremath{\int_{#1}^{#2} #3 \deriv{x}}}				%intégrale dx
\newcommand{\intt}[3]{\ensuremath{\int_{#1}^{#2} #3 \deriv{t}}}				%intégrale dy
\newcommand{\misenforme}{\begin{center} Mis en forme jusqu'ici\\ \line(1,0){400}\\ normalement juste, mais à améliorer depuis ici\end{center}}	%raccourci pour mise en forme
\newcommand*\circled[1]{\tikz[baseline=(char.base)]{
            \node[shape=circle,draw,inner sep=1pt] (char) {#1};}}			%pour entourer un chiffre
\newcommand{\horrule}[1]{\rule{\linewidth}{#1}} 				% Create horizontal rule command with 1 argument of height

\theoremstyle{definition}
\newtheorem{exemp}{Exemple}
\newtheorem{examp}{Example}


%-------------------------------------------
%---Environments----------------------------
%-------------------------------------------
\NewEnviron{boite}[1][0.9]{%
	\begin{center}
		\framecolorbox{red}{white}{%
			\begin{minipage}{#1\textwidth}
 	 			\BODY
			\end{minipage}
		}
	\end{center}
}
\NewEnviron{blackbox}[1][0.9]{%
	\begin{center}
		\framecolorbox{black}{white}{%
			\begin{minipage}{#1\textwidth}
 	 			\BODY
			\end{minipage}
		}
	\end{center}
}
\NewEnviron{exemple}[1][0.8]{%
    \begin{center}
        \framecolorbox{white}{gray!20}{%
            \begin{minipage}{#1\textwidth}
                \begin{exemp}
                    \BODY
                \end{exemp}
            \end{minipage}
        }
    \end{center}
}
\NewEnviron{suiteExemple}[1][0.8]{%
    \begin{center}
        \framecolorbox{white}{gray!20}{%
            \begin{minipage}{#1\textwidth}
                \BODY
            \end{minipage}
        }
    \end{center}
}
\NewEnviron{colExemple}[1][0.8]{%
    \begin{center}
        \framecolorbox{white}{gray!20}{%
            \begin{minipage}{#1\columnwidth}
                \begin{exemp}
                    \BODY
                \end{exemp}
            \end{minipage}
        }
    \end{center}
}
\NewEnviron{example}[1][0.8]{%
    \begin{center}
        \framecolorbox{white}{gray!20}{%
            \begin{minipage}{#1\textwidth}
                \begin{examp}
                    \BODY
                \end{examp}
            \end{minipage}
	}
    \end{center}
}
\NewEnviron{systeq}[1][l]{
			\begin{center}
				$\left\{\begin{array}{#1}
					\BODY
				\end{array}\right.$
			\end{center}
 }





%-------------------------------------------
%---General settings-----------------------
%-------------------------------------------
\renewcommand{\headrulewidth}{1pt}										%ligne au haut de chaque page
\renewcommand{\footrulewidth}{1pt}										%ligne au pied de chaque page
\setstretch{1.6}
\author{Olivier Cloux}


\setstretch{1}
\fancyhead[L]{Sciences de l'information } % formatage en-tête - matière
\fancyhead[C]{Série 10} 	% formatage en-tête - Série
\fancyhead[R]{Olivier Cloux}	% formatage en-tête - Nom
\setmarginsrb{20mm}{20mm}{20mm}{15mm}{-5pt}{15mm}{0pt}{8mm}
\title{	
\normalfont \normalsize 
\textsc{Ecole Polytechnique Fédérale de Lausanne} \\ [25pt] % Your university, school and/or department name(s)
\textsc{Sciences de l'information}\\ [0pt] %Name of the course
\horrule{0.5pt} \\[0.4cm] % Thin top horizontal rule
\huge Série 10 \\ % The assignment title
\horrule{2pt} \\[0.5cm] % Thick bottom horizontal rule
}
\date{}
\renewcommand{\(}{\left(}
\renewcommand{\)}{\right)}
\begin{document}

\maketitle
\setcounter{section}{10}
\subsection{}
Rappelons que $card(X)$ représente le nombre d'éléments dans X (donc $card(A)$ le nombre de symboles de l'alphabet et $card(C)$ le nombre de mots de codes par bloc) ; de plus, n représente la longueur des mots de code, $k$ se calcule par $log_{card(A)}\big(card(C)\big)$, et $d_{min}$ est la plus petite distance de Hamming.
\begin{enumerate}
	\item	\fbox{$n=3$}\\
			$A = \{0,1\} \to card(A) = 2,\ card(C) = 4, \to \log_2(4) =$  \fbox{$k= 2$}\\
			\fbox{$d_{min} = 2$}
			
	\item	\fbox{$n=5$}\\
			$A = \{0,1\} \to card(A) = 2,\ card(C) = 4, \to  \log_2(4) =$ \fbox{$k = 2$}\\
			\fbox{$d_{min} = 3$}
			
	\item	\fbox{$n=6$}\\
			$A = \{0,1,2\} \to card(A) = 3,\ card(C) = 3, \to \log_3(3) =$ \fbox{$k = 1$}\\
			\fbox{$d_{min} = 5$}
\end{enumerate}

\subsection{}
\begin{enumerate}
	\item	\fbox{$A = \{0,1,2,3\}$}

	\item 	$r = \frac{k}{n} = \frac{log_{card(A)}\big(card(C)\big)}{n} = \frac{\log_4(5)}{5} \simeq \frac{1.161}{5} \simeq$ \fbox{$0.232 = r$}

	\item 	\fbox{$d_{min} = 3$} (par exemple entre 01301 et 21310). Cette distance minimale doit répondre à 
			\begin{equation*}
				d_{min}(C) < n(1-r) + 1 = n-k+1 = \text{ borne de Singleton}
			\end{equation*}
			Ainsi, la borne de Singleton est : $n-k+1 = 5-\log_4(5) +1 = 6-\log_4(5) \simeq $ \fbox{4.839}
		
	\item 	Le décodeur d'un code C est capable de corriger tous les effacements de poids $\leq p_1$ si et seulement si $p_1 < d_{min}(C)$. Ainsi, ce code sera capable de corriger tout effacement de poids $d_{min}-1 = 3-2 =$ \fbox{$2 = p_1$}
	
	\item 	La valeur maximale est \fbox{$q_1 = 4$}. En effet, les mots de code $0????,\ 1????,\ 3????$ peuvent être corrigés en respectivement $01301,\ 13012,\ 30122$. $p_2$ est plus grand que $p_1$, car nous savons quels bits sont faux, et dans des cas précis. Il 4 est un chiffre irréaliste pour corriger tous les effacements, car certaines erreurs, même de 3, ne sont plus corrigibles, pour cause de confusion entre deux mots.
	
	\item 	 Le décodeur d'un code C est capable de détecter toutes les erreurs de poids $\leq p_2$ si et seulement si $p_2 < d_{min}(C)$. Ainsi, ce code sera capable de détecter toutes les erreurs de poids $d_{min}-1 = 3-2 =$ \fbox{$2 = p_2$}
	
	\item 	La valeur \fbox{$q_2 = 5$} ; par exemple, si nous recevons le mot de code 00000, nous détectons immédiatement qu'il s'agit d'une erreur, car cela ne correspond à aucun de nos mots de code. Ce n'est pas la même que $p_2$, car comme nous $p_2$ doit fonctionner pour toute erreur. Par exemple, si nous envoyons 01301 mais que par hasard, une erreur parfaite de poids $d_{min}$ arrive, ce mot peut se changer en $21310$. Raison pour laquelle $q_2$ est de 5, alors que $p_2$ est plus petit que $d_{min}$
	
	\item 	Le décodeur d'un code C est capable de corriger toutes les erreurs de poids $\leq p_3$ si et seulement si $p_3 < \frac{d_{min}(C)}{2} = \frac{3}{2} = 1.5$. Le code ne pourra donc corriger que toutes les erreurs de poids \fbox{$p_3 = 1$}.
	
	\item 	Il est imposé que les  mots de code doivent être de longueur 5 ; comme $r = \frac{k}{n} = \frac{k}{5}$, nous devons aussi nous assurer que $k = \log_{card(A)}\big(card(C)\big) = 1$, donc avoir autant de mots de codes que de symboles. De plus, comme il doit corriger toutes les erreurs de poids $\leq 2$, il faut $2 \leq \frac{d_{min}(C)}{2}$, ce qui nous amène à trouver que $d_{min} = 5$. Le code \fbox{$C = \{00000,11111,22222,33333,44444\}$} semble remplir toutes les conditions : La longueur des mots est de 5, il y a autant de mots que de symboles (donc $k=1$ et $r = \frac{1}{5}$). De plus, comme ils sont tous exactement différents, $d_{min} = 5$ et, par le théorème 11.4, il peut corriger les erreurs de poids $\leq 2$
\end{enumerate}

\subsection{}
\begin{enumerate}
	\item 	Nous parlons d'un code binaire de longueur $2^k$; nous avons donc un mot de code $y$, binaire, de longueur $n$; trouver le nombre de mots de codes qui sont à une distance t de y, revient à combiner la distance parmi la longueur. Une combinaison de $t$ parmi $n$ se calcule avec le coefficient binomial $C^t_n = \combi{n}{t} = \frac{n!}{t!(n-t)!}$
	
	\item 	Nous connaissons le binôme de Newton $ \((X+Y)^n = \somme{n}{i=0}\combi{n}{i}X^{n-i}Y^i\)$, ce qui nous montre que $2^n = \somme{n}{i=0}\combi{n}{i}$. Ainsi, selon la donnée, nous aurons $2^k$ mots de code (donc $\somme{n}{i=0}\combi{n}{i}$ mots de codes). Nous savons également, par le dernier exercice, qu'il y a $\combi{x}{e}$ mots d'une longueur donnée qui sont à une distance $e$ de $x$. Comme nous voulons chercher pour tous les $t$ inférieurs à $e$, nous trouvons que le nombre de mots de code dans \fbox{$B(x,e) = \somme{e}{t=0}\combi{n}{t}$}
	
	\item 	Nous savons que la distance de Hamming respecte la symétrie et l'inégalité triangulaire. De plus, comme il est dit que l'intersection des deux ensembles n'est pas vide, cela signifie qu'il existe au moins un mot $y$ valide, tel que $d(x_1,y) \leq e$ et $d(x_2,y) \leq e$. Ensuite, nous savons, pas l'inégalité triangulaire, que $d(x_1,x_2) \leq d(x_1,y) + d(y,x_2) \leq e+e = 2e$. Donc $d(x_1,x_2) \leq 2e$ Notons que nous avons utilisé que $d(x_2,y) = d(y,x_2)$, car cette opération conserve la symétrie.
	
	\item 	Notre code peut corriger toutes les erreurs de poids $\leq e$, ce qui implique que 
			\begin{equation}
				e \leq \frac{d_{min}(C)}{2}
				\label{condi}
			\end{equation}
			De plus, nous venons de prouver que 
			\begin{equation}
				B(x_1,e) \cap B(x_2,e) \neq \emptyset \longrightarrow d(x_1,x_2) \leq 2e
				\label{base}
			\end{equation}
			Par contraposition, nous pouvons donc placer que 
			\begin{equation}
				d(x_1,x_2) > 2e \longrightarrow B(x_1,e) \cap B(x_2,e) = \emptyset
				\label{contra}
			\end{equation}
			Finalement, comme $x_1$ et $x_2$ sont des mots valides distincts de C, il s'ensuit que la distance entre les deux est forcément $\geq d_{min}(C)$ (par définition). Donc $d(x_1,x_2) \geq d_{min}(C)$ et, par \eqref{condi}, $\geq 2e$ ; comme nous l'avons trouvé dans \eqref{contra}, cette condition est suffisante à montrer que l'intersection de ces ensembles est nulle.
			
	\item 	Comme les blocs sont de longueur $n$, il y a un total de $2^n$ mots de codes maximum. Nous avons aussi montré que le nombre de mots de codes dans $B(x,e)$ pour un x donné est de $\somme{e}{t=0}\combi{n}{t}$, et nous savons que le notre code comprend $2^k$ mots de codes. La partie droite de l'inéquation représente donc tous les mots de codes de C multiplié par le nombre de mots \enquote{faux} par mot de code. Cela représente donc l'ensemble des mots qui peuvent être corrigés si besoin (qu'ils soient faux ou non). Mais comme nous savons qu'il n'y a pas de \enquote{doublons} (car $B(x_1,e) \cap B(x_2,e) = \emptyset$), cela signifie qu'il ne peut pas y avoir plus de mots corrigibles que de mots possibles. Il y a donc bien moins de mots qui peuvent être corrigés que de mots possibles, donc $2^n \geq 2^k \somme{e}{t=0}\combi{n}{t}$. 
\end{enumerate}





























\end{document}